% GENERATED by assurance/reports/build_reports.py from the records; do not edit.
\documentclass[11pt,a4paper]{report}
\usepackage{../latex/liondriver-report}
% GENERATED by assurance/reports/build_reports.py from the records; do not edit.
\newcommand{\HbVersion}{0.3}
\newcommand{\HbDate}{2026-11-19}
\newcommand{\PartCount}{18}
\newcommand{\ModeCount}{37}
\newcommand{\RatedCount}{15}
\newcommand{\BlockCount}{7}
\newcommand{\BoxFit}{57.6}
\newcommand{\SgSevenSpf}{37.56}
\newcommand{\SgOneRf}{15.16}
\newcommand{\SgOneLatent}{20.7}
\newcommand{\SgOneSafe}{17.94}
\newcommand{\SpfmBefore}{96.84}
\newcommand{\SpfmNow}{96.96}
\newcommand{\LfmBefore}{88.94}
\newcommand{\LfmNow}{88.54}
\newcommand{\OpenModes}{1}
\newcommand{\AppliedParts}{4}
\newcommand{\FmedaRev}{0.9}
\newcommand{\PeriodicTests}{2}

\LDsetdoc{WP-H-03 Harness Box Component FMEDA}{\HbVersion}
\title{Harness Box Component FMEDA\\[4pt]\large WP-H-03 \S6.5, part list version \HbVersion}
\author{LionDriver assurance\\\small generated from \texttt{assurance/04-hardware/fmeda/harness-box-parts.yaml}}
\date{\HbDate}
\begin{document}
\maketitle
\tableofcontents
\chapter{Introduction}
This report renders the component-level FMEDA of the harness box v3 (element E-05; WP-H-03 \S6.5, part list version \HbVersion{} of \HbDate{}) from its source of record, the hardware-level analysis ISO 26262-5 asks for \cite{iso26262-5}. It takes the box part by part: \PartCount{} parts with \ModeCount{} failure modes, in \BlockCount{} functional blocks, and proves that the parts add up to the box rows of the two goal FMEDAs, for SG-07 (keep the stock pre-collision braking) and SG-01 (lateral control within the envelope).

\paragraph{The box in five lines.}
\begin{enumerate}
\item \PartCount{} parts and \ModeCount{} failure modes, every designator of the v3 drawing and the BOM; \RatedCount{} carry their own rate, the connectors sit inside the wiring row.
\item For stock braking (SG-07) the box puts \SgSevenSpf{}~FIT on single points out of \BoxFit{}~FIT, almost all of it the two DG419LE switches.
\item For steering (SG-01) the box has no single point: \SgOneRf{}~FIT is covered by a mechanism, \SgOneLatent{}~FIT is latent and \SgOneSafe{}~FIT is safe.
\item The part-level view found what the block view hid: the goal FMEDAs split the same switch two ways, and \AppliedParts{} ignition-sense parts were in no goal FMEDA. Both findings are now in the SG-001 FMEDA (revision \FmedaRev{}): the SPFM of SG-01 went from \SpfmBefore{}\,\% to \SpfmNow{}\,\% and its LFM from \LfmBefore{}\,\% to \LfmNow{}\,\%.
\item No periodic test of the box exists in service, so every latent fault stays latent for the box's life; modes whose effect is still to be analysed: \OpenModes{}.
\end{enumerate}
The rates are generic estimates: neither the Siemens failure-rate standard \cite{sn29500} nor a failure-mode handbook is held, and no stress conversion to the mission profile is applied. These numbers show where the box falls short; they are not evidence of compliance. Every number in this report is a macro written from the record (\texttt{numbers.tex}), and every table and finding is computed by the same checks as the record's generated page, so the report cannot fall out of step with the analysis. The record is a draft, AI-assisted (independence I0) and not reviewed; so is this report.
\chapter{The box and its blocks}
Figure~\ref{fig:blocks} shows the functional blocks of the box between the Toyota harness and the comma device; Table~\ref{tab:blocks} gives what each block shall do, its parts and their summed rate.
\begin{figure}[htbp]\centering
\includegraphics[width=\linewidth,height=0.8\textheight,keepaspectratio]{fig-h02-3-9-box-blocks.pdf}
\caption{The functional blocks of the harness box (WP-H-02 Figure 3-9).}\label{fig:blocks}
\end{figure}
{\small
\begin{longtable}{@{}p{1.0cm}p{3.0cm}p{6.4cm}p{2.4cm}p{1.2cm}@{}}
\caption{Functional blocks with their parts and summed rate.}\label{tab:blocks}\\
\toprule \textbf{Block} & \textbf{Name} & \textbf{Shall} & \textbf{Parts} & \textbf{FIT} \\ \midrule\endfirsthead
\toprule \textbf{Block} & \textbf{Name} & \textbf{Shall} & \textbf{Parts} & \textbf{FIT} \\ \midrule\endhead
B1 & Box supply and protection & Supply both switches from the ignition line through D1 (or from comma power on J2), and clamp supply transients above the 16 V stand-off of D4 & D1, D4, J2 & 3 \\ \midrule
B2 & Switch drive & Hold both switch inputs low through 240 $\Omega$, so that the stock path is selected whenever the device does not drive SBU1 & R2, R3 & 0.4 \\ \midrule
B3 & Intercept switch & Join the camera bus to the car bus while undriven (pass-through), and separate them while SBU1 is driven high (intercept), on CAN\_H (U1) and CAN\_L (U2) & U1, U2 & 50 \\ \midrule
B4 & Termination in intercept & Terminate the separated camera bus with 60 $\Omega$ while intercepting (R7, and R8 through the harness loop of pins 13 and 15); R6 serves cars whose plug loops pins 3 and 5, and is unused with Toyota A & R7, R8, R6 & 0.6 \\ \midrule
B5 & Bus clamps & Clamp the CAN0 and CAN2 lines to box VIN and ground & D2, D3 & 2 \\ \midrule
B6 & Ignition sense & Pull SBU2 low through 1 k$\Omega$ while the ignition line is on, so that the device reads the ignition & Q6, R1, R4, R5 & 1.6 \\ \midrule
B7 & Connectors & Carry the car and camera buses, ignition and ground from the harness plug (J1) and the buses, SBU lines and supply to the device cable (J3) & J1, J3 & 0 \\ \midrule
\bottomrule
\end{longtable}
}
\chapter{The FMEDA flow}
Figure~\ref{fig:flow} shows the inputs a detailed component FMEDA needs and the outputs it gives, with the records that hold each here and how far each has got (Tables~\ref{tab:in} and~\ref{tab:out}).
\begin{figure}[htbp]\centering
\includegraphics[width=\linewidth,height=0.75\textheight,keepaspectratio]{fig-h03-6-0-fmeda-flow.pdf}
\caption{The FMEDA flow (WP-H-03 Figure 6-0).}\label{fig:flow}
\end{figure}
{\small
\begin{longtable}{@{}p{4.4cm}p{1.8cm}p{7.8cm}@{}}
\caption{FMEDA inputs.}\label{tab:in}\\
\toprule \textbf{Input} & \textbf{State} & \textbf{Note} \\ \midrule\endfirsthead
\toprule \textbf{Input} & \textbf{State} & \textbf{Note} \\ \midrule\endhead
Field failure data & not needed & Optional. LionDriver has none; proven in use is not claimed (T-06) \\ \midrule
Failure modes and safe state & done & DFMEA modes per element; safe state per safety goal \\ \midrule
Architecture with failure modes per block & done & Block diagram (WP-H-02 §3) and block FMEA (WP-H-03 §4) \\ \midrule
Schematic, BOM, component data & partial & The harness box only (v3 drawing, netlist). The comma device has none (UK-01, UK-02; DQ-065) \\ \midrule
List of automatic diagnostics & done & Safety mechanisms SM-01\ldots{}45 with DC basis \\ \midrule
Prevention measures and periodic tests & partial & DFMEA actions planned; no periodic test of the box in service \\ \midrule
\bottomrule
\end{longtable}
}
{\small
\begin{longtable}{@{}p{4.4cm}p{1.8cm}p{7.8cm}@{}}
\caption{FMEDA outputs.}\label{tab:out}\\
\toprule \textbf{Output} & \textbf{State} & \textbf{Note} \\ \midrule\endfirsthead
\toprule \textbf{Output} & \textbf{State} & \textbf{Note} \\ \midrule\endhead
Failure rate per failure mode and category & partial & Per part and mode, generic rates (no SN 29500 held) \\ \midrule
Useful life & partial & No wear-out part except connector mating cycles, to confirm \\ \midrule
Diagnostic coverage per failure mode & done & From the mechanisms the goal FMEDAs credit \\ \midrule
Coverage of each periodic test & missing & No periodic test of the box exists; coverage 0 in service \\ \midrule
\bottomrule
\end{longtable}
}
\chapter{Failure rate per mode and category}
Figure~\ref{fig:parts} shows every part with its rate and the class of each failure mode for both goals; Table~\ref{tab:cat} sums the rated parts by category and Table~\ref{tab:parts} lists every mode with the goal FMEDA row it rolls up into.
\begin{figure}[htbp]\centering
\includegraphics[width=\linewidth,height=0.85\textheight,keepaspectratio]{fig-h03-6-5-box-parts.pdf}
\caption{The parts, their rates and the class of each mode (WP-H-03 Figure 6-5).}\label{fig:parts}
\end{figure}
\begin{longtable}{@{}p{6.4cm}p{3.4cm}p{3.4cm}@{}}
\caption{Rate of the rated parts per category and goal.}\label{tab:cat}\\
\toprule \textbf{Category} & \textbf{SG-07 braking (FIT)} & \textbf{SG-01 steering (FIT)} \\ \midrule\endfirsthead
\toprule \textbf{Category} & \textbf{SG-07 braking (FIT)} & \textbf{SG-01 steering (FIT)} \\ \midrule\endhead
Single-point & 37.56 & 0 \\ \midrule
Covered by a mechanism (before DC) & 0 & 15.16 \\ \midrule
Latent multiple-point & 3.08 & 20.7 \\ \midrule
Safe & 15.24 & 17.94 \\ \midrule
No effect on this goal & 1.72 & 3.8 \\ \midrule
Effect not yet analysed & 0 & 0 \\ \midrule
\bottomrule
\end{longtable}
{\small
\begin{longtable}{@{}p{0.8cm}p{1.7cm}p{5.4cm}p{1.1cm}p{0.9cm}p{0.9cm}p{2.6cm}@{}}
\caption{Every part and failure mode with its class per goal.}\label{tab:parts}\\
\toprule \textbf{Part} & \textbf{FIT} & \textbf{Failure mode} & \textbf{Share} & \textbf{SG-07} & \textbf{SG-01} & \textbf{Goal FMEDA row} \\ \midrule\endfirsthead
\toprule \textbf{Part} & \textbf{FIT} & \textbf{Failure mode} & \textbf{Share} & \textbf{SG-07} & \textbf{SG-01} & \textbf{Goal FMEDA row} \\ \midrule\endhead
U1 & 25 & Stuck in intercept, or IN input damaged high & 30\% & SPF & safe & sg009 FMD-019, sg001 FMD-021 \\ \midrule
 &  & Stuck in pass-through & 30\% & safe & RF & sg009 FMD-021, sg001 FMD-020 \\ \midrule
 &  & Channel open, increased on-resistance, leakage or latch-up & 40\% & SPF & MPF-L & sg009 FMD-022, sg001 FMD-022 \\ \midrule
U2 & 25 & Stuck in intercept, or IN input damaged high & 30\% & SPF & safe & sg009 FMD-019, sg001 FMD-021 \\ \midrule
 &  & Stuck in pass-through & 30\% & safe & RF & sg009 FMD-021, sg001 FMD-020 \\ \midrule
 &  & Channel open, increased on-resistance, leakage or latch-up & 40\% & SPF & MPF-L & sg009 FMD-022, sg001 FMD-022 \\ \midrule
R2 & 0.2 & Open (IN floats when SBU1 is undriven) & 60\% & MPF-L & safe & sg009 FMD-028, sg001 FMD-024 \\ \midrule
 &  & Short (SBU1 held low) & 40\% & safe & RF & sg009 FMD-029, sg001 FMD-023 \\ \midrule
R3 & 0.2 & Open (IN floats when SBU1 is undriven) & 60\% & MPF-L & safe & sg009 FMD-028, sg001 FMD-024 \\ \midrule
 &  & Short (SBU1 held low) & 40\% & safe & RF & sg009 FMD-029, sg001 FMD-023 \\ \midrule
R7 & 0.2 & Open (camera side left with one 120 $\Omega$ in intercept) & 60\% & MPF-L & none & sg009 FMD-035 \\ \midrule
 &  & Short (camera CAN2\_H to CAN2\_L in intercept) & 40\% & SPF & none & sg009 FMD-036 \\ \midrule
R8 & 0.2 & Open (camera side left with one 120 $\Omega$ in intercept) & 60\% & MPF-L & none & sg009 FMD-035 \\ \midrule
 &  & Short (camera CAN2\_H to CAN2\_L in intercept) & 40\% & SPF & none & sg009 FMD-036 \\ \midrule
R6 & 0.2 & Open & 60\% & none & none & none \\ \midrule
 &  & Short & 40\% & none & none & none \\ \midrule
D1 & 2 & Open (box unpowered without comma power) & 60\% & SPF & safe & sg009 FMD-030 \\ \midrule
 &  & Short (reverse protection lost) & 40\% & MPF-L & none & sg009 FMD-032 \\ \midrule
D4 & 1 & Short (box VIN to ground) & 60\% & SPF & safe & sg009 FMD-031 \\ \midrule
 &  & Open (transient protection lost) & 40\% & MPF-L & none & sg009 FMD-032 \\ \midrule
D2 & 1 & Short or leakage on a CAN line & 30\% & SPF & none & sg009 FMD-033 \\ \midrule
 &  & Open (clamp lost) & 70\% & MPF-L & none & sg009 FMD-034 \\ \midrule
D3 & 1 & Short or leakage on a CAN line & 30\% & SPF & none & sg009 FMD-033 \\ \midrule
 &  & Open (clamp lost) & 70\% & MPF-L & none & sg009 FMD-034 \\ \midrule
Q6 & 1 & Drain-source short (SBU2 held low) & 50\% & none & MPF-L & sg001 FMD-025 \\ \midrule
 &  & Open (SBU2 never pulled low) & 50\% & none & safe & sg001 FMD-026 \\ \midrule
R1 & 0.2 & Open & 60\% & none & safe & sg001 FMD-027 \\ \midrule
 &  & Short & 40\% & safe & safe & sg001 FMD-028 \\ \midrule
R4 & 0.2 & Open & 60\% & none & safe & sg001 FMD-029 \\ \midrule
 &  & Short & 40\% & none & MPF-L & sg001 FMD-030 \\ \midrule
R5 & 0.2 & Open & 60\% & none & MPF-L & sg001 FMD-031 \\ \midrule
 &  & Short & 40\% & none & safe & sg001 FMD-032 \\ \midrule
J1 & in the wiring row & Contact open in the stock path & 30\% & SPF & none & sg009 FMD-023 \\ \midrule
 &  & Plug loose (whole box off the harness) & 70\% & safe & safe & sg009 FMD-024 \\ \midrule
J3 & in the wiring row & SBU1 contact open & 30\% & safe & RF & none \\ \midrule
 &  & Cable loose (device off the box) & 70\% & safe & safe & sg009 FMD-024 \\ \midrule
J2 & no rate & Contact open or short & 100\% & open & open & none \\ \midrule
\bottomrule
\end{longtable}
}
\chapter{Useful life}
Mission: 15 years, about 7,500 operating hours, one power cycle per drive (DQ-067). The rates assume a constant failure rate, which holds only while no part reaches wear-out (Table~\ref{tab:life}). No electrolytic capacitor, relay, battery or fan is on the v3 box, the usual wear-out parts.
{\small
\begin{longtable}{@{}p{3.6cm}p{5.4cm}p{5.0cm}@{}}
\caption{Wear-out per part group.}\label{tab:life}\\
\toprule \textbf{Parts} & \textbf{Wear-out} & \textbf{Limit} \\ \midrule\endfirsthead
\toprule \textbf{Parts} & \textbf{Wear-out} & \textbf{Limit} \\ \midrule\endhead
U1, U2, Q6, D1, D2, D3, D4 & None expected (semiconductors; no electromigration or oxide limit is given in the transcribed data) & Not limiting \\ \midrule
R1, R2, R3, R4, R5, R6, R7, R8 & None expected (thick-film resistors at low power) & Not limiting \\ \midrule
J1, J2, J3 & Contact wear by mating cycles and fretting under vibration & To confirm against the connector datasheets; the box is mated once at installation and again only at service or decommissioning (WP-O-02) \\ \midrule
\bottomrule
\end{longtable}
}
\chapter{Diagnostic coverage per failure mode}
Table~\ref{tab:dc} lists the modes a goal FMEDA credits a safety mechanism for, with its diagnostic coverage. Every other mode of the box has no mechanism: DC 0\,\%.
{\small
\begin{longtable}{@{}p{0.8cm}p{4.6cm}p{1.2cm}p{6.2cm}p{0.9cm}@{}}
\caption{Diagnostic coverage the goal FMEDAs credit.}\label{tab:dc}\\
\toprule \textbf{Part} & \textbf{Failure mode} & \textbf{Goal} & \textbf{Mechanism} & \textbf{DC} \\ \midrule\endfirsthead
\toprule \textbf{Part} & \textbf{Failure mode} & \textbf{Goal} & \textbf{Mechanism} & \textbf{DC} \\ \midrule\endhead
U1 & Stuck in pass-through & SG-01 & FMD-SM-006 & 90\% \\ \midrule
U1 & Channel open, increased on-resistance, leakage or latch-up & SG-01 & FMD-SM-008 & 90\% \\ \midrule
U2 & Stuck in pass-through & SG-01 & FMD-SM-006 & 90\% \\ \midrule
U2 & Channel open, increased on-resistance, leakage or latch-up & SG-01 & FMD-SM-008 & 90\% \\ \midrule
R2 & Short (SBU1 held low) & SG-01 & FMD-SM-006 & 90\% \\ \midrule
R3 & Short (SBU1 held low) & SG-01 & FMD-SM-006 & 90\% \\ \midrule
\bottomrule
\end{longtable}
}
\chapter{Periodic tests}
Table~\ref{tab:tests} lists the \PeriodicTests{} tests that find latent faults of the box, the proof tests of the FMEDA vocabulary. None runs in service.
{\small
\begin{longtable}{@{}p{3.2cm}p{2.4cm}p{3.2cm}p{1.4cm}p{3.8cm}@{}}
\caption{Tests that find latent faults of the box.}\label{tab:tests}\\
\toprule \textbf{Test} & \textbf{When} & \textbf{Finds} & \textbf{In service} & \textbf{Note} \\ \midrule\endfirsthead
\toprule \textbf{Test} & \textbf{When} & \textbf{Finds} & \textbf{In service} & \textbf{Note} \\ \midrule\endhead
Installation check of the stock path (DFM-ACT-008) & Once, at installation (WP-O-01) & U1 channel open, U2 channel open, J1 contact open & 0\% & Runs once, so it is not a periodic test for the latent faults of the 15-year life \\ \midrule
Harness box bench tests VS-HW-22\ldots{}24 & On a sample box, at verification (WP-H-06) & R2/R3 open, R7/R8 open, clamp short split, box VIN lost & 0\% & Type tests on a sample; they find design errors, not latent faults in a fielded box \\ \midrule
\bottomrule
\end{longtable}
}
\chapter{Findings}
\section{Applied to the goal FMEDAs}
The findings of the part-level view are in the SG-001 FMEDA, revision 0.9; the roll-up check now finds the parts and both goal FMEDAs in agreement, row by row.
\begin{itemize}
\item OI-11: The DG419LE split of the SG-001 file aligned with the SG-009 file and this part list (30 \% pass-through, 30 \% intercept, 40 \% channel open).
\item OI-12: Q6, R1, R4 and R5 of the ignition sense carried as FMD-C-009 to FMD-C-012 and classified for SG-01, and a DFMEA row added for the ignition sense reading off while driving (DFM-024).
\end{itemize}
R6 has no effect on either goal in the reference configuration, so it has no goal FMEDA row (Table~\ref{tab:parts}).
\section{SG-01 metrics before and after}
Table~\ref{tab:whatif} gives the SG-01 metrics before the findings were applied, as recorded, and now, computed from the SG-001 FMEDA with the project's FMEDA metrics. The ASIL~C SPFM target of 97\,\% is still missed, by 0.04 points (WP-H-04 OI-6).
\begin{longtable}{@{}p{6.6cm}p{2.0cm}p{2.0cm}p{2.6cm}@{}}
\caption{SG-01 metrics before and after the part-level findings.}\label{tab:whatif}\\
\toprule \textbf{Case} & \textbf{SPFM} & \textbf{LFM} & \textbf{Total (FIT)} \\ \midrule\endfirsthead
\toprule \textbf{Case} & \textbf{SPFM} & \textbf{LFM} & \textbf{Total (FIT)} \\ \midrule\endhead
SG-01 before the findings (fmeda-sg001 rev 0.8) & 96.84\,\% & 88.94\,\% & 255.4 \\ \midrule
SG-01 now (fmeda-sg001 rev 0.9) & 96.96\,\% & 88.54\,\% & 257 \\ \midrule
\bottomrule
\end{longtable}
\section{Modes whose effect is not yet analysed}
\begin{itemize}
\item J2: Contact open or short
\end{itemize}
\appendix
\chapter{Abbreviations}
Written from the project's single list of terms (\texttt{terms.py}), with only the terms this report uses.
\begin{description}\small
\item[AI] artificial intelligence
\item[ASIL] Automotive Safety Integrity Level, A (lowest) to D (highest), from ISO 26262
\item[BOM] bill of materials
\item[CAN] controller area network: the vehicle's data bus
\item[DC] diagnostic coverage: the share of faults a safety mechanism detects
\item[DFM-nn] DFMEA row: failure mode DFM-nnn, cause .Cn (08-analyses/fmea)
\item[DFMEA] design FMEA: failure modes and effects analysis of a design (here the harness box and device)
\item[DQ-nn] decision of the project lead (decision register)
\item[E-nn] element of the item (WP-C-01 §3.1), E-01\ldots{}E-06
\item[FIT] failures in time: failures per 10$^9$ hours
\item[FMD-nn] FMEDA row: component FMD-C-nnn, failure mode FMD-nnn (08-analyses/fmea)
\item[FMEA] failure mode and effects analysis
\item[FMEDA] failure modes, effects and diagnostic analysis (hardware metrics)
\item[I0–I3] independence level of a review, I0 (none, e.g. the author or an AI) to I3 (independent organisation)
\item[ISO] International Organization for Standardization
\item[LFM] latent fault metric: share of the multiple-point failure rate that is not latent (ISO 26262-5)
\item[MPF-L] latent multiple-point fault: violates a safety goal only together with another fault, and nothing detects it
\item[OI-nn] open item of the document it appears in
\item[RF] residual fault: the part of a fault a safety mechanism does not cover (ISO 26262-5)
\item[SBU] sideband-use pins of a USB-C connector (in the harness, the intercept drive, orientation and ignition lines)
\item[SG-nn] safety goal (HARA §6)
\item[SM-nn] safety mechanism (TSC §6)
\item[SN] SN 29500, the Siemens standard of reference failure rates for electronic components
\item[SPF] single-point fault: a fault no safety mechanism covers that violates a safety goal on its own (ISO 26262-5)
\item[SPFM] single-point fault metric: share of the safety-related failure rate not single-point or residual (ISO 26262-5)
\item[T-nn] code-derived detection or reaction time (WP-S-04 §4); in WP-M-01 and WP-S-06, a tailoring decision T-01\ldots{}T-08 (WP-M-01 §5)
\item[UK-nn] unknown that WP-H-02 must obtain (UK-01\ldots{}UK-09, WP-H-02 §5)
\item[VIN] vehicle identification number; in the harness box and its simulations, the box supply voltage
\item[VS-nn] verification case: VS-SI system integration, VS-HW hardware, VS-SWQ software, VS-FI fault injection, VS-SQ system qualification
\item[WP-nn] work product, e.g. WP-C-04; the letter is its area (M management, C concept, S system, H hardware, W software, V validation, P processes, A analyses, O operation, K safety case)
\end{description}
\begin{thebibliography}{9}
\bibitem{iso26262-5} ISO 26262-5:2018, Road vehicles: Functional safety. Part 5: Product development at the hardware level. ISO, Geneva.
\bibitem{sn29500} SN 29500, Failure rates of components: Expected values (series). Siemens AG, Munich.
\end{thebibliography}
\end{document}
